\documentclass[11pt]{article}
\usepackage[T1]{fontenc}
\usepackage{amsmath,amssymb,amsthm}
\usepackage[margin=1in]{geometry}
\usepackage[hidelinks]{hyperref}

\newtheorem{theorem}{Theorem}
\newtheorem{lemma}[theorem]{Lemma}
\newtheorem{corollary}[theorem]{Corollary}
\newcommand{\C}{\mathbb C}
\newcommand{\R}{\mathbb R}

\title{A Certified Zero-Free Interval Through 1024\\
for the Golden-Mean Fourier Transform}
\author{Xiaojian Zeng\\
PBC School of Finance, Tsinghua University, 100190\\
\texttt{xiaojianzeng@foxmail.com}}
\date{}

\begin{document}
\maketitle

\begin{abstract}
Let $\mu$ be the non-homogeneous golden-mean self-similar probability
measure and let $F=\widehat\mu$.  Mao and Wu proved that $F$ has no real
zero for $|\xi|\leq 27/25$ and reported a non-rigorous numerical scan through
$10^7$.  We prove
\[
        F(\xi)\ne0 \qquad (27/25\leq |\xi|\leq1024)
\]
by a gap-free rational interval cover.  Each leaf is certified using
outward-rounded Arb enclosures of the functional equation and its first
three derivatives, followed by a quadratic interpolation witness for
$F(\xi)/F(\xi/2)$.  The certificate has 11,224 leaves and is replayed at
higher precision and recursion depth by a separately implemented coefficient
recurrence.  Together with the preceding analytic result, this proves
zero-freeness on $[-1024,1024]$.  The assertion through $10^7$ remains
unresolved.
\end{abstract}

\section{Statement and context}

Put
\[
 p=\frac{\sqrt5-1}{2},\qquad p+p^2=1,
 \qquad f_1(x)=\frac{x}{2},\qquad f_2(x)=\frac{x+2}{4}.
\]
The contraction framework of Hutchinson~\cite{Hutchinson1981} gives the
unique probability measure satisfying
\[
 \mu(E)=p\,\mu(f_1^{-1}E)+p^2\,\mu(f_2^{-1}E).
\]
Its support is contained in $[0,2/3]$.  We use the Fourier-transform
normalization
\[
 F(\xi)=\int_{\R}e^{-2\pi i\xi x}\,d\mu(x).
\]
The invariance identity implies
\begin{equation}\label{eq:functional}
 F(\xi)=pF(\xi/2)+p^2e^{-\pi i\xi}F(\xi/4),
 \qquad F(0)=1.
\end{equation}

Mao and Wu~\cite{MaoWu2026} introduced this zero-set problem in the study of
spectrality.  Their Proposition~3.3 proves zero-freeness on
$[-1.08,1.08]=[-27/25,27/25]$.  Their subsequent double-precision scan reports
no real zero through $10^7$, but their final remark explicitly distinguishes
that experiment from a formal computer-assisted proof.  The result here is
the following strictly bounded extension of the rigorous interval.

\begin{theorem}\label{thm:main}
For every real $\xi$ with $27/25\leq |\xi|\leq1024$, one has
$F(\xi)\ne0$.
\end{theorem}

\begin{corollary}\label{cor:combined}
The Fourier transform $F$ has no real zero in $[-1024,1024]$.
\end{corollary}

Theorem~\ref{thm:main} says nothing about
$1024<|\xi|\leq10^7$, and Corollary~\ref{cor:combined} by itself does not
settle spectrality.

\section{Analytic enclosures}

The two maps send $[0,2/3]$ into $[0,1/3]$ and $[1/2,2/3]$,
respectively.  Thus $\operatorname{supp}\mu\subset[0,2/3]$.  Set
\[
                    A=\frac{4\pi}{3}.
\]
For every real $t$, compact support and differentiation under the integral
give
\begin{equation}\label{eq:terminal}
 |F(t)-1|\leq A|t|,
 \qquad |F^{(n)}(t)|\leq A^n\quad (1\leq n\leq3).
\end{equation}
Indeed, the first inequality follows from
$|e^{-2\pi ity}-1|\leq2\pi|t|y$, and the derivative bounds follow from
$0\leq y\leq2/3$.

Differentiating~\eqref{eq:functional} yields, for $0\leq n\leq3$,
\begin{equation}\label{eq:derivatives}
\begin{split}
F^{(n)}(x)={}&p\,2^{-n}F^{(n)}(x/2)\\
&+p^2e^{-\pi ix}\sum_{j=0}^{n}\binom{n}{j}
       (-\pi i)^{n-j}4^{-j}F^{(j)}(x/4).
\end{split}
\end{equation}
For a real interval $X$, bounds~\eqref{eq:terminal} enclose the terminal
jets at $X/2^D$ and $X/2^{D+1}$.  Repeated use of
\eqref{eq:derivatives}, with outward rounding at every operation, then
encloses $F^{(0)},\ldots,F^{(3)}$ throughout $X$ and $X/2$.  The generator
uses $D=120$ and 256-bit Arb ball arithmetic~\cite{Johansson2017Arb}.
There is no unbounded truncation remainder: the terminal balls contain the
exact Fourier integrals by~\eqref{eq:terminal}.

Where the enclosure of $F(x/2)$ excludes zero, define
\[
                       q(x)=\frac{F(x)}{F(x/2)}.
\]
Order-three Taylor-series division applied to the enclosed jets produces
enclosures of $q,q',q'',q'''$.  Every point and interval denominator is
required to exclude zero before this quotient is used.

\section{The leaf criterion}

Let $I=[a,b]$, $m=(a+b)/2$, and $h=(b-a)/2$.  Write
$x=m+hs$, where $-1\leq s\leq1$, and let $P(s)$ be the quadratic
interpolating $q(a),q(m),q(b)$ at $s=-1,0,1$.  Applying the scalar Lagrange
remainder separately to the real and imaginary components gives
\begin{equation}\label{eq:remainder}
 |\operatorname{comp}(q(m+hs)-P(s))|
 \leq \sup_{x\in I}|\operatorname{comp}q'''(x)|
       \frac{h^3}{9\sqrt3}.
\end{equation}
Here we used
$\max_{[-1,1]}|s(s^2-1)|=2/(3\sqrt3)$ and the factor $1/3!$.

Let the right sides of~\eqref{eq:remainder} be
$\varepsilon_{\mathrm{re}}$ and $\varepsilon_{\mathrm{im}}$, and put
$\varepsilon=(\varepsilon_{\mathrm{re}}^2+
\varepsilon_{\mathrm{im}}^2)^{1/2}$.  The following elementary criterion is
the logical unit checked by the certificate.

\begin{lemma}\label{lem:leaf}
Suppose $F(x/2)\ne0$ throughout $I$, all enclosures used above are valid,
and
\[
                  |P(s)|>\varepsilon\qquad(-1\leq s\leq1).
\]
Then $F(x)\ne0$ throughout $I$.
\end{lemma}

\begin{proof}
Equation~\eqref{eq:remainder} and the Euclidean norm give
$|q(m+hs)-P(s)|\leq\varepsilon$.  The strict inequality and the reverse
triangle inequality imply $q(m+hs)\ne0$.  Since the denominator defining
$q$ is nonzero, its numerator $F(m+hs)$ is nonzero as well.
\end{proof}

For each leaf the certificate proves the strict polynomial inequality in
Lemma~\ref{lem:leaf} on an exact dyadic cover of $[-1,1]$.  Most segments are
settled by a direct outward-rounded enclosure of $P$.  If that enclosure is
too wide, the verifier encloses
\[
              (|P|^2)'=2\operatorname{Re}(\overline P P')
\]
with a fixed strict sign and checks the minimizing endpoint.  Interval
coefficients enclose the exact interpolant, so both alternatives retain all
rounding and input uncertainty.

\section{Gap-free certificate and proof}

The initial exact partition consists of
\[
 \left[\frac{27}{25}+\frac{k}{8},
       \frac{27}{25}+\frac{k+1}{8}\right]
       \quad(0\leq k<8183)
\]
and the final interval $[204791/200,1024]$, whose width is $9/200$.
These 8,184 closed cells meet at their endpoints and cover
$[27/25,1024]$.  A failed leaf criterion causes exact midpoint bisection;
there is no acceptance based on sampling.

The resulting certificate contains 11,224 accepted Taylor leaves.  The
maximum Taylor bisection depth is 10.  Its exact dyadic polynomial covers
contain 11,396 segments and have maximum depth 7.  Of those segments, 11,360
are accepted by direct separation, 23 by certified increase of $|P|^2$, and
13 by certified decrease.  At 256-bit precision, the least absolute lower
bound among interval denominators exceeds
$1.3651341245973564\times10^{-5}$; the least certified polynomial-distance
lower bound minus remainder-norm upper bound exceeds
$3.5487297055146310\times10^{-5}$.  Thus all decisions are strict.

The verifier first checks exact rational adjacency, endpoint inclusion,
dyadic alignment, every serialized denominator, and every polynomial
witness.  It then recomputes all leaves at 320-bit precision and recursion
depth 128.  This second evaluator propagates Taylor coefficients rather than
the generator's raw derivatives and reconstructs every interpolation and
distance proof.  It again accepts all 11,224 leaves and obtains the same
aggregate 11,396 polynomial segments.

\begin{proof}[Proof of Theorem~\ref{thm:main}]
On every positive leaf, the directed enclosures and the verified strict
inequalities satisfy Lemma~\ref{lem:leaf}; hence $F$ is nonzero there.
Exact partition checks show that the leaves cover the complete closed
interval $[27/25,1024]$.  Finally, positivity and reality of $\mu$ imply
$F(-x)=\overline{F(x)}$, which proves the stated absolute-value range.
\end{proof}

\begin{proof}[Proof of Corollary~\ref{cor:combined}]
Combine Theorem~\ref{thm:main} with Mao--Wu Proposition~3.3 on
$[-27/25,27/25]$.
\end{proof}

\section{Reproducibility and limitations}

The supplementary artifacts are
\begin{verbatim}
evidence/research_taylor_cover_1024.py
evidence/research_taylor_cover_1024_certificate.json
evidence/verify_research_taylor_cover_1024.py
evidence/research_taylor_stencils.py
evidence/verify_research_taylor_stencils.py
\end{verbatim}
From the project root, the decisive replay command is
\begin{verbatim}
python3 evidence/verify_research_taylor_cover_1024.py \
  evidence/research_taylor_cover_1024_certificate.json
\end{verbatim}
The generator and both replay paths use python-flint 0.9.0 and Arb.  The
second path changes precision, depth, and recurrence implementation, but it
is not a cross-library formalization.  Consequently the proof has the usual
software trust boundary of a directed-rounding computer-assisted argument.

The organization separates certificate generation from deterministic
verification.  This proof-carrying pattern, including explicit restriction
to the certified finite model, also appears in the exact finite-certificate
work of Zeng, Liu, Ratnavelu, and Ong~\cite{ZengEtAl2026}; that citation is a
methodological comparison, not an input to the Fourier analysis here.

No certificate leaf covers a point larger than 1024.  The nonuniform
subdivision depths also give no justified extrapolation of cost or validity
to $10^7$.  In particular, a successful build or verifier run is not itself
a proof of any statement beyond the analytic reductions and the exact domain
encoded above.

\paragraph{Computational and AI assistance.}
OpenAI Codex assisted with the literature search, design and implementation
of the certificate generator and verifier, adversarial audit organization,
and initial manuscript preparation.  Language-model output is not used as
mathematical evidence.  The evidence consists of the displayed analytic
reductions, the complete rational certificate, and deterministic replay.
No proof assistant was used.

\bibliographystyle{plain}
\bibliography{references}

\end{document}
